\section{Coding 为什么像 GPU：加速器而不是应用}

把 Coding 模型类比为 GPU，是本期最有解释力的比喻。GPU 本身不是一个消费应用，但它让训练和推理成为可能；Coding 模型也不只是 IDE 插件，而是让研究、工程、数据处理、实验和产品迭代加速的基础设施。

\subsection{三种加速路径}

\noindent\begin{tabularx}{\textwidth}{p{0.22\textwidth}p{0.34\textwidth}X}
\toprule
加速路径 & 发生在哪里 & 影响 \\
\midrule
个人生产力放大 & 顶尖研究员、程序员、架构师 & 1\% 人才被放大 10--50 倍，想法到实验的周期压缩。 \\
组织研发吞吐 & feature、数据 pipeline、多模态实验 & 原本数周的迭代缩到数天，产品发布频率上升。 \\
AI 加速 AI & 模型帮助写训练代码、评估脚本、调试实验 & AI 研究本身被 AI 工具加速，形成递归反馈。 \\
\bottomrule
\end{tabularx}

\begin{knowledgebox}{为什么代码是数字世界的杠杆}
自然语言表达问题，代码表达解决方案。只要任务发生在数字系统里，代码就能连接数据、工具、环境、权限和执行。模型掌握 coding，等于掌握了数字世界的操作杆。
\end{knowledgebox}

\subsection{本章小结}

Coding 是加速器，不是单一应用。它放大个人、组织和 AI 研究自身，因此可能成为 AGI 第二幕的关键基础设施。

